



\begin{table*}[t]
\centering
\caption{\textbf{[RQ1] Effectiveness of \sys.} This table shows the attack effectiveness of \sys across four different datasets, three different models, and three different recomputation methods. The KV cache system are conducted under default settings with a cache ratio $\delta = 0.3$ and a recomputation ratio $\rho = 0.1$. \textit{Acc} is the model's task accuracy reported to measure the impact of position-independent KV reuse on model utility. \textit{U-ASR} and \textit{T-ASR} is reported as the attack performance metrics. }
\label{tab:rq1_main}

\footnotesize
\setlength{\tabcolsep}{0pt}

\begin{tabular*}{\textwidth}{@{\extracolsep{\fill}} l ccc ccc ccc ccc }
\toprule
\multirow{2}{*}{\textbf{Method}} & \multicolumn{3}{c}{\textbf{HotpotQA}} & \multicolumn{3}{c}{\textbf{SQuAD}} & \multicolumn{3}{c}{\textbf{MedQA}} & \multicolumn{3}{c}{\textbf{PubMedQA}} \\
\cmidrule(lr){2-4} \cmidrule(lr){5-7} \cmidrule(lr){8-10} \cmidrule(lr){11-13}
 & Acc & U-ASR & T-ASR & Acc & U-ASR & T-ASR & Acc & U-ASR & T-ASR & Acc & U-ASR & T-ASR \\
\midrule

\rowcolor{sectionbg}
\multicolumn{13}{l}{\textbf{Llama-3.1-8B}} \\
Full        & 0.73 & --   & --   & 0.76 & --   & --   & 0.71 & --   & --   & 0.92 & -- & -- \\
Vanilla     & 0.58 & 1.00 & 1.00 & 0.60 & 0.98 & 0.96 & 0.55 & 0.98 & 0.90 & 0.74 & 0.93 & 0.92 \\
Random      & 0.63 & 0.95 & 0.94 & 0.69 & 0.98 & 0.89 & 0.63 & 0.93 & 0.87 & 0.79 & 0.92 & 0.89 \\
EPIC        & 0.66 & 0.87 & 0.85 & 0.72 & 0.89 & 0.75 & 0.70 & 0.96 & 0.82 & 0.87 & 0.96 & 0.91 \\
CacheBlend  & 0.68 & 0.92 & 0.89 & 0.72 & 0.94 & 0.85 & 0.68 & 0.88 & 0.80 & 0.89 & 0.92 & 0.86 \\

\rowcolor{sectionbg}
\multicolumn{13}{l}{\textbf{Ministral-8B}} \\
Full        & 0.70 & --   & --   & 0.67 & --   & --   & 0.68 & --   & --   & 0.94 & -- & -- \\
Vanilla     & 0.57 & 1.00 & 1.00 & 0.55 & 0.98 & 0.93 & 0.54 & 1.00 & 1.00 & 0.71 & 0.97 & 0.94 \\
Random      & 0.61 & 0.97 & 0.97 & 0.62 & 0.89 & 0.87 & 0.63 & 1.00 & 1.00 & 0.75 & 0.91 & 0.88 \\
EPIC        & 0.67 & 0.96 & 0.95 & 0.65 & 0.88 & 0.78 & 0.63 & 1.00 & 1.00 & 0.84 & 0.92 & 0.87 \\
CacheBlend  & 0.66 & 0.94 & 0.94 & 0.66 & 0.93 & 0.87 & 0.65 & 1.00 & 1.00 & 0.85 & 0.89 & 0.85 \\

\rowcolor{sectionbg}
\multicolumn{13}{l}{\textbf{Qwen3-8B}} \\
Full        & 0.82 & --   & --   & 0.84 & --   & --   & 0.73 & --   & --   & 0.95 & -- & -- \\
Vanilla     & 0.61 & 0.96 & 0.94 & 0.68 & 1.00 & 1.00 & 0.57 & 0.97 & 0.96 & 0.75 & 1.00 & 1.00 \\
Random      & 0.72 & 0.91 & 0.85 & 0.74 & 0.95 & 0.90 & 0.66 & 0.86 & 0.83 & 0.85 & 0.93 & 0.89 \\
EPIC        & 0.75 & 0.88 & 0.87 & 0.75 & 0.92 & 0.86 & 0.71 & 0.93 & 0.84 & 0.92 & 0.95 & 0.93 \\
CacheBlend  & 0.80 & 0.91 & 0.89 & 0.82 & 0.93 & 0.87 & 0.70 & 0.92 & 0.87 & 0.92 & 1.00 & 0.91 \\

\bottomrule
\end{tabular*}
    \vspace{-0.1in}
\end{table*}
\subsection{RQ1: Effectiveness}\label{sec:rq1}

\paragraph{Setup.}
RQ1 investigate the effectiveness of the \sys across different models, datasets, and recomputation methods.
% 
Experiments are conducted on four QA benchmarks: HotpotQA~\cite{yang2018hotpotqa}, SQuAD (v1.1 and v2.0)~\cite{rajpurkar2016squad,rajpurkar2018know}, MedQA~\cite{jin2021disease}, and PubMedQA~\cite{jin2019pubmedqa}. 
% 
Regarding hyperparameters, we set the cache ratio $\delta = 0.3$, the recomputation ratio $\rho = 0.1$, and the chunk size $L_{\mathrm{chunk}} = 32$ as the default settings~\cite{yao2025cacheblend, hu2025epic}. We evaluate the effectiveness of \sys without using KV cache (\textit{Full}), with full KV cache reuse (\textit{Vanilla}), and with three recomputation methods (\textit{Random}, \textit{EPIC}, and \textit{CacheBlend}).
% 
We use white-box setting and report Acc for model performance, and U-ASR and T-ASR for the effectiveness of \sys.


\paragraph{Results.} 
Table~\ref{tab:rq1_main} shows that \sys demonstrates strong attack performance across all evaluated scenarios.
% 
Specifically, on the Llama-3.1-8B model, \sys achieves an average T-ASR of 89\% across the four datasets, even when countering the four distinct recomputation methods. While advanced recomputation strategies like EPIC and CacheBlend result in a slight reduction in ASR compared to the Vanilla setting, \sys still maintains  effectiveness. This trend is further supported by the results on Qwen3-8B, where the attack frequently achieves near-perfect success rates (e.g., 100\% T-ASR on PubMedQA), proving that \sys remains robust despite the slight drop introduced by partial recomputation.

\paragraph{Conclusion.}
In conclusion, \sys is a highly effective, model-agnostic attack that hijacks LLM generation across diverse domains. Our findings indicate that standard partial recomputation mechanisms, at current ratios, are insufficient to mitigate the adversarial cache optimized by \sys.